% Part: normal-modal-logic
% Chapter: filtrations
% Section: preliminaries

\documentclass[../../../include/open-logic-section]{subfiles}

\begin{document}

\olfileid{nml}{fil}{pre}

\olsection{Prasyarat}

Filtrasi ngidini kita netepake katerputusan sistem logika modal kanthi
nuduhake yen sistem-sistem kasebut nduweni \emph{sipat model
  winates}, yaiku saben !!{formula} kang bener (salah) ing sawijining
model uga bener (salah) ing sawijining model \emph{winates}. Filtrasi
ditepakake relatif marang himpunan !!{formula}s kang katutup tumrap
subformula.

\begin{defn}\ollabel{defn:modallyclosed}
  Himpunan $\Gamma$ saka !!{formula}s iku \emph{katutup tumrap
  subformula} yen ngemot saben subformula saka !!a{formula} ing
  $\Gamma$. Sabanjure, $\Gamma$ iku \emph{katutup modal} yen
  katutup tumrap subformula lan, saliyane iku, $!A \in \Gamma$ njalari
  $\Box!A, \Diamond!A \in \Gamma$.
\end{defn}

Upamane, kanggo !!a{formula}~$!A$, himpunan kabeh
sub-!!{formula}s-ne katutup tumrap sub-!!{formula}s. Nalika kita
netepake filtrasi sawijining model liwat himpunan sub-!!{formula}s
saka~$!A$, filtrasi iku bakal nduweni sipat kang dikarepake: filtrasi
ndadekake $!A$ bener (salah) yen lan mung yen model asline mangkono.

Himpunan jagad saka filtrasi~$\mModel{M}$ liwat~$\Gamma$ ditepakake
minangka himpunan kabeh kelas ekuivalensi saka relasi ekuivalensi ing
ngisor iki.

\begin{defn}
Upamane $\mModel{M} =\tuple{W, R, V}$ lan $\Gamma$ katutup tumrap
sub-!!{formula}s. Tetepna relasi $\equiv$ ing $W$ kang nggayutake
rong jagad yen !!{formula}s saka $\Gamma$ kang bener ing loro-lorone
padha, yaiku:
\[
u \equiv v \quad \text{yen lan mung yen } \quad
\forall !A \in \Gamma : \mSat{M}{!A}[u] \Leftrightarrow \mSat{M}{!A}[v].
\]
Kelas ekuivalensi~$[w]_\equiv$ saka jagad~$w$, utawa cekake $[w]$,
yaiku himpunan kabeh jagad kang $\equiv$-ekuivalen karo~$w$:
\[
[w] = \Setabs{v}{v \equiv w}.
\]
\end{defn}

\begin{prop}
  Kanggo $\mModel{M}$ lan $\Gamma$, $\equiv$ kang ditepakake ing
  ndhuwur iku relasi ekuivalensi, yaiku refleksif, simetris, lan
  transitif.
\end{prop}

\begin{proof}
  Relasi $\equiv$ iku refleksif, amarga $w$ ndadekake persis
  !!{formula}s saka~$\Gamma$ kang padha bener kaya awake dhewe. Relasi
  iki simetris, amarga yen $u$ ndadekake !!{formula}s saka~$\Gamma$
  kang padha bener kaya $v$, kosok baline uga bener kanggo $v$ lan~$u$.
  Relasi iki uga transitif, amarga yen $u$ ndadekake !!{formula}s
  saka~$\Gamma$ kang padha bener kaya~$v$, lan $v$ kaya $w$, mula
  $u$ ndadekake !!{formula}s saka~$\Gamma$ kang padha bener kaya~$w$.
\end{proof}

Relasi $\equiv$, kaya saben relasi ekuivalensi, mbagi~$W$ dadi
\emph{kelas-kelas partisi}, yaiku himpunan bagean saka~$W$ kang
padha-padha ora tumpang tindih lan bebarengan nutupi kabeh~$W$.
Saben $w \in W$ iku !!a{element} saka salah siji kelas partisi,
yaiku~$[w]$, amarga $w \equiv w$. Dadi kelas-kelas $[w]$ nutupi
kabeh~$W$. Kelas-kelas iku padha-padha ora tumpang tindih, amarga yen
$u \in [w]$ lan $u \in [v]$, mula $u \equiv w$ lan $u \equiv v$;
miturut simetri lan transitivitas, $w \equiv v$, mula $[w] = [v]$.

\end{document}
